% Options for packages loaded elsewhere
\PassOptionsToPackage{unicode}{hyperref}
\PassOptionsToPackage{hyphens}{url}
\PassOptionsToPackage{dvipsnames,svgnames,x11names}{xcolor}
%
\documentclass[
  11pt,
  a4paper,
]{article}

\usepackage{amsmath,amssymb}
\usepackage{iftex}
\ifPDFTeX
  \usepackage[T1]{fontenc}
  \usepackage[utf8]{inputenc}
  \usepackage{textcomp} % provide euro and other symbols
\else % if luatex or xetex
  \usepackage{unicode-math}
  \defaultfontfeatures{Scale=MatchLowercase}
  \defaultfontfeatures[\rmfamily]{Ligatures=TeX,Scale=1}
\fi
\usepackage[]{libertinus}
\ifPDFTeX\else  
    % xetex/luatex font selection
    \setmonofont[Scale=0.92]{Latin Modern Mono}
\fi
% Use upquote if available, for straight quotes in verbatim environments
\IfFileExists{upquote.sty}{\usepackage{upquote}}{}
\IfFileExists{microtype.sty}{% use microtype if available
  \usepackage[]{microtype}
  \UseMicrotypeSet[protrusion]{basicmath} % disable protrusion for tt fonts
}{}
\makeatletter
\@ifundefined{KOMAClassName}{% if non-KOMA class
  \IfFileExists{parskip.sty}{%
    \usepackage{parskip}
  }{% else
    \setlength{\parindent}{0pt}
    \setlength{\parskip}{6pt plus 2pt minus 1pt}}
}{% if KOMA class
  \KOMAoptions{parskip=half}}
\makeatother
\usepackage{xcolor}
\usepackage[a4paper,textheight=24cm,textwidth=15.5cm]{geometry}
\usepackage{svg}
\setlength{\emergencystretch}{3em} % prevent overfull lines
\setcounter{secnumdepth}{5}
% Make \paragraph and \subparagraph free-standing
\makeatletter
\ifx\paragraph\undefined\else
  \let\oldparagraph\paragraph
  \renewcommand{\paragraph}{
    \@ifstar
      \xxxParagraphStar
      \xxxParagraphNoStar
  }
  \newcommand{\xxxParagraphStar}[1]{\oldparagraph*{#1}\mbox{}}
  \newcommand{\xxxParagraphNoStar}[1]{\oldparagraph{#1}\mbox{}}
\fi
\ifx\subparagraph\undefined\else
  \let\oldsubparagraph\subparagraph
  \renewcommand{\subparagraph}{
    \@ifstar
      \xxxSubParagraphStar
      \xxxSubParagraphNoStar
  }
  \newcommand{\xxxSubParagraphStar}[1]{\oldsubparagraph*{#1}\mbox{}}
  \newcommand{\xxxSubParagraphNoStar}[1]{\oldsubparagraph{#1}\mbox{}}
\fi
\makeatother


\providecommand{\tightlist}{%
  \setlength{\itemsep}{0pt}\setlength{\parskip}{0pt}}\usepackage{longtable,booktabs,array}
\usepackage{calc} % for calculating minipage widths
% Correct order of tables after \paragraph or \subparagraph
\usepackage{etoolbox}
\makeatletter
\patchcmd\longtable{\par}{\if@noskipsec\mbox{}\fi\par}{}{}
\makeatother
% Allow footnotes in longtable head/foot
\IfFileExists{footnotehyper.sty}{\usepackage{footnotehyper}}{\usepackage{footnote}}
\makesavenoteenv{longtable}
\usepackage{graphicx}
\makeatletter
\newsavebox\pandoc@box
\newcommand*\pandocbounded[1]{% scales image to fit in text height/width
  \sbox\pandoc@box{#1}%
  \Gscale@div\@tempa{\textheight}{\dimexpr\ht\pandoc@box+\dp\pandoc@box\relax}%
  \Gscale@div\@tempb{\linewidth}{\wd\pandoc@box}%
  \ifdim\@tempb\p@<\@tempa\p@\let\@tempa\@tempb\fi% select the smaller of both
  \ifdim\@tempa\p@<\p@\scalebox{\@tempa}{\usebox\pandoc@box}%
  \else\usebox{\pandoc@box}%
  \fi%
}
% Set default figure placement to htbp
\def\fps@figure{htbp}
\makeatother

\usepackage{tikz}
\usepackage{xcolor}
\usepackage{tabularx}
\usepackage{orcidlink}
\usepackage{lineno}
\usepackage{libertinus}
\usepackage{draftwatermark}

\makeatletter
\@ifpackageloaded{caption}{}{\usepackage{caption}}
\AtBeginDocument{%
\ifdefined\contentsname
  \renewcommand*\contentsname{Table of contents}
\else
  \newcommand\contentsname{Table of contents}
\fi
\ifdefined\listfigurename
  \renewcommand*\listfigurename{List of Figures}
\else
  \newcommand\listfigurename{List of Figures}
\fi
\ifdefined\listtablename
  \renewcommand*\listtablename{List of Tables}
\else
  \newcommand\listtablename{List of Tables}
\fi
\ifdefined\figurename
  \renewcommand*\figurename{Figure}
\else
  \newcommand\figurename{Figure}
\fi
\ifdefined\tablename
  \renewcommand*\tablename{Table}
\else
  \newcommand\tablename{Table}
\fi
}
\@ifpackageloaded{float}{}{\usepackage{float}}
\floatstyle{ruled}
\@ifundefined{c@chapter}{\newfloat{codelisting}{h}{lop}}{\newfloat{codelisting}{h}{lop}[chapter]}
\floatname{codelisting}{Listing}
\newcommand*\listoflistings{\listof{codelisting}{List of Listings}}
\makeatother
\makeatletter
\makeatother
\makeatletter
\@ifpackageloaded{caption}{}{\usepackage{caption}}
\@ifpackageloaded{subcaption}{}{\usepackage{subcaption}}
\makeatother
\makeatletter
\@ifpackageloaded{algorithm}{}{\usepackage{algorithm}}
\makeatother
\makeatletter
\@ifpackageloaded{algpseudocode}{}{\usepackage{algpseudocode}}
\makeatother
\makeatletter
\@ifpackageloaded{caption}{}{\usepackage{caption}}
\makeatother
    \makeatletter
    \newcommand\fs@nocaption{
      \def\@fs@cfont{\bfseries}
      \let\@fs@capt\floatc@plain
      \def\@fs@pre{}%
      \def\@fs@post{\kern2pt\hrule}%
      \def\@fs@mid{\hrule\kern2pt}%
      \let\@fs@iftopcapt\iftrue}
    \makeatother
  
\captionsetup[algorithm]{justification=raggedright}

\usepackage{bookmark}

\IfFileExists{xurl.sty}{\usepackage{xurl}}{} % add URL line breaks if available
\urlstyle{same} % disable monospaced font for URLs
\hypersetup{
  pdftitle={Causal inference of post-transcriptional regulation timelines from long-read sequencing in Arabidopsis thaliana},
  pdfauthor={Rubén Martos; Christophe Ambroise; Guillem Rigaill},
  pdfkeywords={Arabidopsis thaliana, causal discovery, causal
inference, chloroplast, DAG, do-calculus, EM-algorithm, genomics, graphical
models, imputation, intervention transcription},
  colorlinks=true,
  linkcolor={blue},
  filecolor={Maroon},
  citecolor={Blue},
  urlcolor={Blue},
  pdfcreator={LaTeX via pandoc}}


\title{Causal inference of post-transcriptional regulation timelines
from long-read sequencing in Arabidopsis thaliana}
\author{Rubén Martos \and Christophe Ambroise \and Guillem Rigaill}
\date{2026-03-28}

\begin{document}
\definecolor{computo-blue}{HTML}{034E79}

\begin{tikzpicture}[remember picture,overlay]
\fill[computo-blue]
  (current page.north west) -- (current page.north east) --
  ([yshift=-5cm]current page.north east|-current page.north east) --
  ([yshift=-5cm]current page.north west|-current page.north west) -- cycle;
\node[anchor=north west, xshift=.75cm,
  yshift=-.75cm] at (current page.north west) {\includegraphics[height=3cm]{logo_text_white}};
\node[font=\sffamily\bfseries\color{white},anchor=north west, xshift=.75cm,
  yshift=-4.25cm] at (current page.north
  west) {\fontsize{10}{12}\selectfont ISSN 2824-7795};
\node[font=\sffamily\bfseries\color{white},anchor=west,
  xshift=4.25cm,yshift=-2.75cm] at (current page.north west)
  {\begin{minipage}{15cm}
    \fontsize{25}{30}\selectfont
    Causal inference of post-transcriptional regulation timelines from
    long-read sequencing in Arabidopsis thaliana
    \vspace{.5cm}
    \\
    \fontsize{15}{18}\selectfont
    
  \end{minipage}};
\end{tikzpicture}

\vspace*{2.5cm}
\begin{center}
          Rubén
Martos~\orcidlink{0000-0002-1463-5088}\footnote{Corresponding author: \href{mailto:ruben.martosprieto@univ-evry.fr}{ruben.martosprieto@univ-evry.fr}}\quad
             LaMME, Université d'Évry Paris-Saclay\\
                 Christophe
Ambroise~\orcidlink{0000-0002-8148-0346}\quad
             LaMME, Université d'Évry Paris-Saclay\\
                 Guillem Rigaill~\orcidlink{0000-0002-7176-7511}\quad
             IPS2, LaMME, Université Paris-Saclay / CNRS / INRAE /
Université d'Évry\\
           
  \bigskip
  
  Date published: 2026-03-28 \quad Last modified: 2026-03-28
\end{center}
      
\bigskip
\begin{abstract}
We propose a novel framework for reconstructing the chronology of
genetic regulation using causal inference based on Pearl's theory. The
approach proceeds in three main stages: causal discovery, causal
inference, and chronology construction. We apply it to the \emph{ndhB}
and \emph{ndhD} genes of the chloroplast in \emph{Arabidopsis thaliana},
generating four alternative maturation timeline models per gene, each
derived from a different causal discovery algorithm (HC, PC, LiNGAM, or
NOTEARS). Two methodological challenges are addressed: the presence of
missing data, handled via an EM algorithm that jointly imputes missing
values and estimates the Bayesian network, and the selection of the
\(\ell_1\)-regularization parameter in NOTEARS, for which we introduce a
stability selection strategy. The resulting causal models consistently
outperform reference chronologies in terms of both reliability and model
fit. Moreover, by combining causal reasoning with domain expertise, the
framework enables the formulation of testable hypotheses and the design
of targeted experimental interventions grounded in theoretical
predictions.
\end{abstract}

\noindent%
{\it Keywords:} Arabidopsis thaliana, causal discovery, causal
inference, chloroplast, DAG, do-calculus, EM-algorithm, genomics, graphical
models, imputation, intervention transcription
\vfill

\linenumbers
\DraftwatermarkOptions{stamp=true, colorspec=0.9, text={\rm submitted}}

\floatname{algorithm}{Algorithm}

\renewcommand*\contentsname{Contents}
{
\hypersetup{linkcolor=}
\setcounter{tocdepth}{3}
\tableofcontents
}

\section{Causal inference of post-transcriptional regulation timelines
from long-read sequencing in Arabidopsis thaliana}\label{section}

2026-03-28

\emph{We present a framework for reconstructing the chronology of
chloroplast gene regulation in Arabidopsis thaliana using causal
inference, addressing missing data challenges, and generating testable
models for ndhB and ndhD genes.}

\href{https://github.com/rmartosprieto/computo-chloroDAG/actions/workflows/build.yml}{\pandocbounded{\includesvg[keepaspectratio]{README_files/mediabag/badge.svg}}}
\href{http://creativecommons.org/licenses/by/4.0/}{\pandocbounded{\includegraphics[keepaspectratio]{README_files/mediabag/80x15.png}}}

\subsubsection{Authors}\label{authors}

\begin{itemize}
\tightlist
\item
  \href{}{Rubén Martos} (Université d'Évry Paris-Saclay)
\item
  \href{}{Christophe Ambroise} (Université d'Évry Paris-Saclay)
\item
  \href{}{Guillem Rigaill} (Université Paris-Saclay / CNRS / INRAE /
  Université d'Évry)
\end{itemize}

\subsubsection{Abstract}\label{abstract}

We propose a novel framework for reconstructing the chronology of
genetic regulation using causal inference based on Pearl's theory. The
approach proceeds in three main stages: causal discovery, causal
inference, and chronology construction. We apply it to the \emph{ndhB}
and \emph{ndhD} genes of the chloroplast in \emph{Arabidopsis thaliana},
generating four alternative maturation timeline models per gene, each
derived from a different causal discovery algorithm (HC, PC, LiNGAM, or
NOTEARS). Two methodological challenges are addressed: the presence of
missing data, handled via an EM algorithm that jointly imputes missing
values and estimates the Bayesian network, and the selection of the
\(\ell_1\)-regularization parameter in NOTEARS, for which we introduce a
stability selection strategy. The resulting causal models consistently
outperform reference chronologies in terms of both reliability and model
fit. Moreover, by combining causal reasoning with domain expertise, the
framework enables the formulation of testable hypotheses and the design
of targeted experimental interventions grounded in theoretical
predictions.




\end{document}
